\section{Related Work: Pre-Execution Commitment versus Adjacent Frames}
\label{sec:related}

We organize related work into ten substantive clusters plus a methodological reflexivity cluster and a brief summary. Each cluster occupies a slightly different temporal or structural slice of the agent pipeline. We position our contribution by contrast: the pre-execution responsibility commitment that this paper foregrounds is not the same object as task type, ambiguous information need, missing input, model confidence, behavioural variance, post-hoc failure mode, output-aggregation choice, judge bias, decomposed evaluation rubric, or annotator disagreement, even though every prior cluster shares some commitment with our framing.

\subsection{Task-aware delegation, routing, and cost-aware model selection}
\label{sec:related:task_aware}

The closest neighbour is task-aware delegation, which derives task-conditioned capability and coordination-risk signals from preference data and uses them to drive routing, auditor selection, and rationale disclosure~\cite{taskaware2026}. Its loop --- task typing $\to$ capability profile $\to$ coordination-risk cue $\to$ routing decision $\to$ accountability log --- shares our commitment to making delegation visible and auditable, to surfacing reliability cues to users, and to logging accountability after the fact. A complementary line of routing-and-cost work designs systems that pick a model per query under cost and quality trade-offs~\cite{ong2024routellm, routerbench2024, flyswatcannon2023, cqr2024costefficient, chen2023reducecost}. Routing benchmarks such as RouterBench~\cite{routerbench2024} stress-test these systems across heterogeneous tasks, and meta-modeling approaches learn the routing policy from preference data~\cite{ong2024routellm} or from cost/performance prediction labels~\cite{flyswatcannon2023}. The shared commitment is that typed delegation is the right granularity. The difference is the typed object: routing systems condition on a task cluster $c$ and ask ``which agent is reliable on this task type?'', while we condition on a responsibility vector $r$ over a closed dimension set and ask ``what responsibility structure has the agent projected from this delegation, and should it be clarified before commitment?'' Routing returns a pointer; projection returns a structured commitment that can be partial, refused, or explicitly claimed.

\subsection{Ambiguity identification and clarification}
\label{sec:related:ambiguity}

A long line of work studies whether language models can recognize and clarify ambiguous user queries. ClariQ~\cite{aliannejadi2020clariq} studies clarification in open-domain information-seeking dialogue with a benchmark of clarifying-question turns; CLAMBER~\cite{clamber2024} more recently organizes ambiguity along epistemic-misalignment, linguistic-ambiguity, and aleatoric-output axes and tests whether language models generate useful clarifying questions. Open-domain QA work treats ambiguity as a target signal: AmbigQA constructs disambiguation pairs for ambiguous open-domain questions~\cite{min2020ambigqa}, and learning-to-ask methods rank candidate clarification questions by expected information value~\cite{rao2018learn2ask}. Recent work also asks when to clarify and when to answer: selective-clarification methods refuse only on the questions where clarification reduces error, leaving the rest answered~\cite{lee2023selectiveambig, cole2023selectivelyambig}. The shared commitment with our protocol is that clarification is an admissible mid-pipeline action; the difference is the object of clarification. Ambiguity benchmarks treat the user's information need as the ambiguous object (``when did he land on the moon?'' --- unclear referent). We treat the agent's responsibility commitment as the ambiguous object (``improve this paper draft.'' --- unclear whether grammar polish, conceptual reconstruction, novelty audit, or all of these are being delegated).

\subsection{Underspecification in long-horizon tasks}
\label{sec:related:underspec}

Long-horizon underspecification work generates controllably underspecified tasks by removing information along Goal, Constraint, Input, and Context axes, then validates the resulting variants behaviourally as outcome-critical, divergent, or benign~\cite{lhaw2024}. It also operationalizes clarification efficiency as a gain-per-question metric, internalizing the intuition that clarification is costly. We adopt the cost-sensitivity insight directly: the protocol's clarification trigger (Equation~(\ref{eq:trigger})) is conditioned on expected loss reduction net of clarification cost, not on ambiguity flags alone. The difference is what is missing: long-horizon underspecification studies missing task information; we study ambiguous responsibility boundaries on a delegation that may already specify the input fully.

\subsection{Uncertainty communication and tool-use calibration}
\label{sec:related:calibration}

How an agent communicates uncertainty is a separate measurement track. Calibration work elicits confidence scores from language models and asks whether expressed confidence matches accuracy: verbalized-confidence prompts~\cite{tian2023justask, lin2022teaching}, internal-token-probability methods~\cite{kadavath2022mostlyknow}, and empirical evaluations across elicitation strategies~\cite{xiong2024canllms}. A separate strand asks whether \emph{users} accurately estimate model knowledge~\cite{steyvers2025whatpeoplethinkknow}, and a recent metacognition study examines whether supervised fine-tuning improves uncertainty communication and whether such improvement transfers across tasks~\cite{metacog2024}. Tool-use calibration work disaggregates calibration further by tool type, observing that evidence tools (web search, retrieval) can induce overconfidence by surfacing noisy material as authoritative, while verification tools (code interpreters) ground reasoning through deterministic feedback~\cite{confidencedichotomy2024}. We treat uncertainty as one cross-cutting responsibility dimension (RX.1) rather than as the main object, and we treat tool outputs as evidence whose reliability $\kappa(e)$ depends on source type (Equation~(\ref{eq:kappa})). The difference is scope: calibration work asks whether confidence matches correctness on the answer; we ask whether the agent disclosed uncertainty about the responsibility structure it claimed to assume.

\subsection{Behavioural variance and reproducibility}
\label{sec:related:variance}

Empirical studies show that an identical prompt and an identical task can produce substantially different analytical conclusions across models, prompting strategies, and temperatures~\cite{sameprompt2024}. Closely related, work on consistency in software-engineering agents shows that consistency correlates with accuracy --- but also that consistency can amplify wrong interpretations, where many failures are ``consistent wrong interpretation'' rather than random error~\cite{consistency2024}. Even nominally deterministic settings exhibit non-determinism in LLM outputs across re-runs~\cite{atkinson2024nondeterminism}. The shared commitment is that variance and interpretation matter as much as raw accuracy. The difference is when in the pipeline we measure: behavioural-variance studies measure outcome and trajectory variance \emph{after} execution; we measure responsibility-projection divergence \emph{before} execution begins. Section~\ref{sec:results:exp1b} shows that pre-execution variance, captured as cross-family projection mismatch, is family-attributable rather than stochastic on the pilot dataset.

\subsection{Multi-agent collaboration frameworks and failure taxonomies}
\label{sec:related:mast}

Multi-agent LLM frameworks specify how agents collaborate to perform consequential work. MetaGPT structures agents as software roles with a meta-programming layer~\cite{hong2024metagpt}; AutoGen builds extensible multi-agent conversations~\cite{wu2024autogen}; CAMEL studies role-playing communicative agents for cooperative task solving~\cite{li2023camel}; ChatDev simulates a software-development team~\cite{qian2024chatdev}; Generative Agents simulates social behaviour~\cite{park2023generative}; and multi-agent debate improves factuality through reasoned disagreement~\cite{du2024debate}. Recent work introduces post-hoc failure taxonomies for these systems, identifying categories such as system-design issues, inter-agent misalignment, and task-verification failures~\cite{mast2025}. These taxonomies are derived from execution traces and label what went wrong after the fact. Among the named modes, ``disobey task specification,'' ``disobey role specification,'' and ``fail to ask for clarification'' overlap directly with our pre-execution governance concerns. The difference is temporal: collaboration frameworks coordinate already-typed roles, and failure taxonomies diagnose completed traces, while we govern the responsibility commitment that may or may not produce the trace in the first place. The frameworks and taxonomies are complementary to our protocol: a multi-agent framework could absorb the projection step as a pre-conversation contract, and the post-hoc taxonomy labels can serve as settlement evidence.

\subsection{Aggregation and the selection bottleneck}
\label{sec:related:selection}

When diverse multi-agent teams produce candidate outputs, the value of diversity depends on the quality of the aggregator: with a strong selector, diverse generation pays off; with a weak selector, it does not~\cite{whenagree2024}. A scaling-laws perspective asks whether more LLM calls reliably improve compound inference systems; the answer is dim-dependent and saturating~\cite{chen2024scaling}. This line asks ``given candidates, which output should be selected?'' We instead ask ``before any output is generated, which responsibility should be projected and committed to?'' These are dual questions on opposite sides of generation. We include a generate-then-select baseline in the experimental design (deferred to the main run; Section~\ref{sec:results:limitations}) so that the cost of generating wrong-scope outputs and selecting among them can be compared directly to the cost of governing scope before generation.

\subsection{LLM-as-judge bias and panel design}
\label{sec:related:judge_bias}

Using LLMs as judges for open-ended outputs introduces bias modes that the panel design must control. Position bias, verbosity bias, and self-enhancement bias are documented baselines~\cite{zheng2023mtbench, wang2023notfair}. LLM judges are inconsistent across re-runs and biased in dataset-specific ways~\cite{stureborg2024inconsistent}. Self-preference / narcissistic-evaluator effects appear when judges identify and favor their own outputs~\cite{panickssery2024selfpref, koo2024narcissistic, wataoka2024selfpreference}. Length-controlled evaluators correct one specific bias (verbosity) but do not eliminate the rest~\cite{dubois2024lengthcontrolled}. Replacing single judges with juries of diverse models is one mitigation~\cite{verga2024juries}; specialized evaluator models (Prometheus 2~\cite{kim2024prometheus2}, G-Eval with GPT-4 prompts~\cite{liu2023geval}) are another. Our 12-judge tier-stratified Anthropic-excluded panel (Section~\ref{sec:results:exp2}) directly absorbs these mitigations: family exclusion controls self-preference, tier-stratification controls capability-tier confound, and panel size targets reliability. The difference from this cluster is the object being judged: judge-bias work asks ``does the judge correctly compare two outputs?'' while we ask ``does the executor's responsibility projection match a panel-derived ground truth before execution?''

\subsection{Decomposed and rubric-based evaluation}
\label{sec:related:decomposed}

Rather than scoring an output as a single quality number, decomposed evaluation maps the output to per-skill or per-requirement components. FLASK decomposes into a fine-grained skill set and scores each skill independently~\cite{ye2024flask}; InFoBench decomposes complex instructions into per-requirement satisfaction~\cite{qin2024infobench}; G-Eval scores along chained explanation steps~\cite{liu2023geval}; Prometheus 2 trains a dedicated evaluator language model on rubric criteria~\cite{kim2024prometheus2}. Our closed responsibility taxonomy $J_{v1.1}$ shares the per-component structure: weights $r_j$ are predicted per dimension, and settlement is computed per dimension. The difference is the decomposition target: FLASK / InFoBench / G-Eval decompose the \emph{output} (skills demonstrated, requirements satisfied, reasoning steps); we decompose the \emph{pre-execution commitment} (responsibility dimensions projected). Per-component aggregation is the shared paradigm; the per-component object differs.

\subsection{Annotator disagreement, human-label variation, and crowdworker reliability}
\label{sec:related:annotator}

Ground-truth annotation tracks have known reliability limits. Plank's framework reframes human-label variation as signal rather than noise~\cite{plank2022humanlabel}, calling for distributional ground truths over single labels. Krippendorff's $\alpha$ formalizes inter-annotator reliability as the canonical measure for interval-level coding data~\cite{krippendorff2011alpha, hayes2007standard}, and our $\alpha$ reporting (Section~\ref{sec:results:dataset:alpha}) follows this convention. The crowdworker-reliability literature documents that even careful Mechanical-Turk evaluations of generated text can systematically differ from expert evaluations~\cite{karpinska2021turkperils, clark2021allthatsgold}, and recent observational work shows that crowdworkers increasingly use LLMs themselves to produce annotations~\cite{veselovsky2023crowdai}, with related work finding systematic variation in human-respondent data quality across online sample-provider platforms~\cite{frenkenberg2025aiagentprev}. Our Stage 1 R1.7 closure (Section~\ref{sec:results:limitations:specifiability}) confirms the crowdworker-reliability concern empirically: the Prolific pass rate against quality screening was 33\%, and the dominant rejection pattern was AI-assisted submissions or task-misunderstanding. The difference between this cluster's framing and ours is what the human-anchor scalability constraint is taken to be: this cluster treats it as a recruitment / training / aggregation problem, while our boundary-condition finding (Section~\ref{sec:discussion:specifiability}) identifies anchor specifiability and domain-expertise gating as binding even after recruitment and aggregation are well-handled.

\subsection{Reasoning, tool use, and self-refinement}
\label{sec:related:reasoning}

A separate line of work increases output quality by changing the executor's reasoning style: chain-of-thought prompting~\cite{wei2022cot}, self-consistency over multiple reasoning paths~\cite{wang2023selfconsistency}, action-reasoning interleaving~\cite{yao2023react}, tool-augmented generation~\cite{schick2023toolformer}, and iterative self-feedback~\cite{madaan2023selfrefine}. Software-engineering work surveys how these techniques are applied to coding agents~\cite{hou2024llm4se}. The shared commitment is that the executor's intermediate steps matter; the difference is what is being intermediated. These techniques alter how the executor reasons \emph{within} an already-committed task; we alter what the executor commits to \emph{before} reasoning begins. The two layers compose: an executor that has projected its responsibility could subsequently reason via chain-of-thought or self-refine, and our protocol does not prescribe a particular reasoning style.

\subsection{Benchmarking the benchmarks}
\label{sec:related:benchmarking}

A reflexive thread asks what counts as a credible NLP benchmark~\cite{bowman2022benchmarking}. The recommendation is to fix construct validity, sample size, and analysis pre-registration before declaring an effect. Our pilot adopts $r^*_{\text{median}}$ pre-registered active-set rules (Section~\ref{sec:experiments}), pre-registered hypothesis tests with paired bootstrap CIs and Krippendorff's $\alpha$ (Section~\ref{sec:experiments:stats}), and explicit reporting of negative results (Section~\ref{sec:results:exp2}). This cluster is methodological and shapes how we report; it is not a competing object.

\subsection{Summary of differentiation}
\label{sec:related:summary}

Across these ten clusters, the recurring contrast is the temporal locus of the typed object. Task-aware delegation and routing fix the typed object as a task cluster; ambiguity benchmarks fix it as a query gap; long-horizon underspecification fixes it as missing input; calibration work fixes it as confidence; behavioural-variance work fixes it as outcome variance; MAS frameworks and failure taxonomies fix roles or post-hoc failure modes; aggregation work fixes the candidate output; LLM-as-judge work fixes the comparison protocol; decomposed evaluation fixes the per-component output rubric; and annotator-disagreement work fixes the human-label distribution. Pre-execution responsibility commitment is the typed object that sits between natural-language delegation and execution, and the rest of the paper develops the closed-taxonomy formulation, the protocol, and the pilot evidence for measuring and acting on it.
